\documentclass[11pt,a4paper]{article}

% ---------------------------------------------------------------- paquets
\usepackage[utf8]{inputenc}
\usepackage[T1]{fontenc}
\usepackage[french]{babel}
\usepackage{lmodern}
\usepackage[margin=2.2cm]{geometry}
\usepackage{xcolor}
\usepackage{hyperref}
\usepackage{booktabs}
\usepackage{longtable}
\usepackage{array}
\usepackage{enumitem}
\usepackage{listings}
\usepackage[most]{tcolorbox}
\usepackage{tikz}
\usetikzlibrary{positioning,arrows.meta,shapes.geometric,fit,calc}

\hypersetup{colorlinks=true,linkcolor=blue!50!black,urlcolor=blue!60!black}
\setlength{\parskip}{0.45em}
\setlength{\parindent}{0pt}
\renewcommand{\arraystretch}{1.25}

\definecolor{bleunuit}{HTML}{1F4E79}
\definecolor{donnees}{HTML}{DBEAFE}
\definecolor{stockage}{HTML}{FEF9C3}
\definecolor{traitement}{HTML}{E0F2E9}
\definecolor{llm}{HTML}{FDE9D9}
\definecolor{interface}{HTML}{EDE9FE}

\newtcolorbox{analogie}{colback=blue!4,colframe=bleunuit,fonttitle=\bfseries,
  title={Analogie},left=6pt,right=6pt,top=4pt,bottom=4pt}
\newtcolorbox{exemple}[1][]{colback=green!4,colframe=green!40!black,fonttitle=\bfseries,
  title={Exemple#1},left=6pt,right=6pt,top=4pt,bottom=4pt}
\newtcolorbox{attention}{colback=orange!6,colframe=orange!70!black,fonttitle=\bfseries,
  title={À retenir},left=6pt,right=6pt,top=4pt,bottom=4pt}

\lstset{basicstyle=\ttfamily\small,breaklines=true,columns=fullflexible,
  frame=single,backgroundcolor=\color{gray!6},rulecolor=\color{gray!40},
  xleftmargin=4pt,xrightmargin=4pt}

\tikzset{
  boite/.style={draw=gray!70,rounded corners=3pt,align=center,font=\small,
    minimum height=1.1cm,text width=2.6cm,inner sep=4pt},
  fleche/.style={-{Stealth[length=2.5mm]},thick,gray!80!black},
  pointille/.style={-{Stealth[length=2.5mm]},thick,dashed,gray!70}
}

\newcommand{\code}[1]{\texttt{#1}}

\title{\textbf{Research Paper Explorer}\\[0.4em]
\large Le projet expliqué simplement, pas à pas}
\author{Projet du module \emph{Generative AI} --- ENSIAS}
\date{}

\begin{document}
\maketitle

\begin{abstract}
Ce document explique, sans prérequis, ce que fait le projet \emph{Research Paper Explorer},
les notions d'intelligence artificielle sur lesquelles il repose, les données utilisées,
l'architecture du code, les raisons de chaque choix technique, les protections contre les
erreurs de l'IA, les limites du travail et les questions qu'un jury peut poser. Chaque notion
est accompagnée d'une analogie de la vie courante et, quand c'est utile, d'un exemple
concret tiré du projet.
\end{abstract}

\tableofcontents
\newpage

% ======================================================================
\section{Le projet en bref}
% ======================================================================

\subsection{Le problème}

Quand on commence un sujet de recherche en informatique, on se heurte à deux difficultés.

\begin{enumerate}
  \item \textbf{Le volume.} Rien que dans la catégorie arXiv \og Calcul et langage \fg{}
  (\code{cs.CL}, le traitement automatique des langues), plus de 81\,000 articles ont été
  publiés depuis janvier 2023. Personne ne peut tout lire, ni même tout parcourir.
  \item \textbf{La fiabilité des assistants IA.} Si l'on demande à un chatbot généraliste
  \og quels articles parlent de tel sujet ? \fg{}, il peut répondre avec assurance\ldots{}
  en inventant des titres et des références qui n'existent pas. C'est ce qu'on appelle une
  \emph{hallucination}.
\end{enumerate}

\subsection{L'objectif}

\emph{Research Paper Explorer} est un assistant de recherche bibliographique qui combine
un moteur de recherche et un modèle de langage, avec une règle stricte : \textbf{chaque
affirmation de la réponse doit citer un article arXiv réel, présent dans les documents
réellement retrouvés}, et cette citation est vérifiée automatiquement par le programme.
Si les documents ne suffisent pas, l'assistant répond \og Je ne sais pas \fg{} plutôt que
d'inventer.

\subsection{Ce que l'utilisatrice peut faire}

L'application web (Streamlit) s'ouvre directement sur la page \textbf{Chat}.

\begin{itemize}[leftmargin=1.2em]
  \item \textbf{Poser une question} en langage naturel, sur tout le corpus ou sur un seul
  article, et obtenir une réponse en français ou en anglais avec des citations
  \code{[arXiv:ID]} cliquables vers \url{https://arxiv.org}.
  \item \textbf{Voir le contexte récupéré} : les articles exacts que l'IA a eus sous les
  yeux pour répondre.
  \item \textbf{Comparer les articles cités} : un tableau (méthode, jeux de données,
  résultats, limites, pertinence), des onglets de lecture et une synthèse comparative,
  exportables en Markdown ou CSV.
  \item \textbf{Générer un rapport PDF} de la réponse et de la comparaison.
  \item \textbf{Retrouver ses conversations} : elles sont enregistrées et listées dans la
  barre latérale ; on peut en ouvrir n'importe laquelle, la télécharger ou la supprimer.
  \item \textbf{Rechercher} des articles sur la page \textbf{Recherche}, avec des filtres
  (années, catégories, conférence, code disponible, article accepté), puis envoyer un
  résultat au Chat avec le bouton \og Interroger avec le RAG \fg{}.
\end{itemize}

\begin{exemple}[ : un parcours typique]
L'utilisatrice tape : \emph{How can retrieval reduce hallucinations in LLMs?}

\begin{enumerate}[nosep]
  \item Le programme cherche les 5 articles les plus pertinents parmi 20\,000.
  \item Le modèle de langage rédige une réponse en s'appuyant uniquement sur ces 5 articles,
  par exemple : \og Combiner un générateur avec des récupérateurs comme BM25 ou DPR réduit
  les hallucinations [arXiv:2506.12483]. \fg{}
  \item Le programme vérifie que \code{2506.12483} fait bien partie des 5 articles. Sinon,
  la citation est supprimée.
  \item Sous la réponse s'affiche la comparaison des articles cités, puis le bouton
  \og Générer le rapport PDF \fg{}.
\end{enumerate}
\end{exemple}

\subsection{Le projet en chiffres}

\begin{center}
\begin{tabular}{ll}
\toprule
Élément & Valeur \\
\midrule
Lignes lues dans le fichier arXiv & 3\,215\,335 \\
Articles \code{cs.CL} publiés depuis 2023 & 81\,503 \\
Articles conservés dans le corpus & 20\,000 (1\,000 en mode de développement) \\
Code & 38 modules Python, environ 4\,200 lignes \\
Tests automatisés & 114 \\
\bottomrule
\end{tabular}
\end{center}

% ======================================================================
\section{Les notions de base, expliquées simplement}
% ======================================================================

\subsection{Intelligence artificielle générative et LLM}

Un \textbf{LLM} (\emph{Large Language Model}, grand modèle de langage) est un programme
entraîné sur d'énormes quantités de texte pour prédire le mot suivant d'une phrase. À force
de prédire le mot suivant, il sait rédiger des paragraphes entiers, résumer, traduire ou
répondre à des questions. C'est de l'IA \emph{générative} : elle produit du texte nouveau.

\begin{analogie}
Un LLM ressemble à une personne qui a lu une bibliothèque entière mais qui répond de
mémoire, sans pouvoir rouvrir les livres. Elle s'exprime très bien, mais quand sa mémoire
est floue, elle peut \og combler les trous \fg{} avec quelque chose de plausible mais faux.
\end{analogie}

Dans le projet, le LLM par défaut est \code{qwen2.5:3b} (3 milliards de paramètres), exécuté
localement grâce au logiciel \textbf{Ollama}. On peut le remplacer par un modèle hébergé
(Groq, Gemini, OpenAI) en modifiant seulement le fichier \code{.env}.

\textbf{Le prompt} est le texte envoyé au LLM : les consignes et la question. Tous les
prompts du projet sont rangés dans un seul fichier, \code{src/rpe/llm/prompts.py}.

\subsection{Les hallucinations}

Une hallucination est une réponse fausse présentée avec assurance : une référence inventée,
un chiffre qui n'apparaît nulle part, un résultat attribué au mauvais article. Le projet est
conçu pour les empêcher ou les détecter (section~\ref{sec:gardefous}).

\subsection{Les embeddings (plongements sémantiques)}

Un \textbf{embedding} transforme un texte en une liste de nombres (un \emph{vecteur}). Dans
le projet, chaque article (\og titre. résumé \fg{}) devient un vecteur de 384 nombres. Deux
textes qui parlent de la même chose ont des vecteurs proches, même s'ils n'emploient pas les
mêmes mots.

\begin{analogie}
Imaginez une immense carte où chaque article est un point. Les articles sur la traduction
automatique sont regroupés dans un quartier, ceux sur la reconnaissance de la parole dans un
autre. Un embedding, c'est l'adresse GPS d'un texte sur cette carte. Chercher des articles
proches d'une question, c'est chercher les points les plus proches de l'adresse de la
question.
\end{analogie}

La proximité se mesure avec la \textbf{similarité cosinus} : plus elle est proche de 1, plus
les deux textes ont un sens voisin.

\begin{exemple}
\og reduce hallucinations \fg{} et \og improve factuality \fg{} n'ont aucun mot en commun,
mais leurs embeddings sont proches : une recherche par embeddings les rapproche.
\end{exemple}

\subsection{La base vectorielle}

Une \textbf{base vectorielle} (ici \textbf{ChromaDB}) stocke les vecteurs de tous les
articles et sait retrouver très vite les plus proches d'un vecteur donné. Elle stocke aussi
des informations à côté de chaque vecteur (année, conférence, code disponible\ldots), ce qui
permet de filtrer.

\subsection{BM25 : la recherche par mots-clés}

\textbf{BM25} est une méthode classique de recherche (celle des moteurs de recherche depuis
les années 1990). Elle donne un score à chaque document selon les mots de la requête qu'il
contient, en accordant plus d'importance aux mots \emph{rares} (un mot présent partout,
comme \og model \fg{}, ne dit pas grand-chose).

\begin{analogie}
BM25, c'est l'index à la fin d'un livre : on cherche le mot exact \og SQuAD \fg{} et on
trouve les pages qui le contiennent. Les embeddings, c'est un bibliothécaire qui comprend
ce que vous voulez dire. Le premier est imbattable pour les noms propres et les sigles, le
second pour les idées.
\end{analogie}

\subsection{La recherche hybride et la fusion RRF}

Le projet utilise \textbf{les deux} moteurs. Chacun propose 100 candidats, puis on les
fusionne avec la \textbf{Reciprocal Rank Fusion} (RRF) :
\[
\text{score}(d) = \sum_{\text{moteurs}} \frac{1}{60 + \text{rang de } d \text{ dans ce moteur}}
\]

\begin{exemple}[ : calcul de la fusion]
L'article A est 1\textsuperscript{er} pour BM25 et 3\textsuperscript{e} pour les
embeddings : $\frac{1}{61}+\frac{1}{63}\approx 0{,}0323$.
L'article B est 1\textsuperscript{er} pour les embeddings mais absent de BM25 :
$\frac{1}{61}\approx 0{,}0164$.
A passe devant B : un article bien classé par les deux moteurs est récompensé.
\end{exemple}

Après la fusion, on garde les \textbf{5 meilleurs articles} pour la réponse du Chat
(paramètre \code{rag\_top\_k} de \code{config.yaml}).

\subsection{Le RAG (Retrieval-Augmented Generation)}

Le \textbf{RAG}, ou \og génération augmentée par la recherche \fg{}, consiste à
\textbf{chercher d'abord} les documents pertinents, puis à demander au LLM de
\textbf{répondre uniquement à partir de ces documents}, en les citant.

\begin{analogie}
Sans RAG, l'étudiant passe un examen à livre fermé : il répond de mémoire. Avec le RAG, il
passe un examen à livre ouvert, mais on ne lui donne que 5 pages précises, et il doit
indiquer pour chaque phrase la page d'où elle vient. Le correcteur (le programme) vérifie
ensuite que chaque numéro de page cité existe bien dans les 5 pages fournies.
\end{analogie}

\subsection{Les sorties structurées (JSON et Pydantic)}

Parfois, on ne veut pas un paragraphe mais des \textbf{champs précis} : la méthode, les jeux
de données, les résultats\ldots{} On demande alors au LLM un objet \textbf{JSON} (un format
texte structuré), et la bibliothèque \textbf{Pydantic} vérifie qu'il a la bonne forme. Si le
JSON est invalide, le programme renvoie l'erreur au LLM et lui demande de corriger.

\begin{analogie}
C'est un formulaire administratif : au lieu d'une lettre libre, on impose des cases à
remplir. Pydantic est l'agent au guichet qui refuse le formulaire s'il manque une case ou si
une date est écrite n'importe comment.
\end{analogie}

\subsection{Notions du cahier des charges qui ne sont plus dans la version livrée}

Le cahier des charges initial prévoyait aussi des modules d'analyse. Ils ont été développés
et testés pendant le projet, puis \textbf{retirés de la version livrée} pour la centrer sur
la recherche et le RAG. Ils restent dans l'historique Git. Ces notions peuvent tout de même
être demandées par le jury.

\paragraph{Le graphe de connaissances et le GraphRAG.}
Un \textbf{graphe de connaissances} relie des entités par des relations : un article
\emph{utilise} le jeu de données SQuAD, \emph{traite} la tâche \og traduction
automatique \fg{}, \emph{est écrit par} tel auteur. Le \textbf{GraphRAG} interroge ce graphe
pour répondre à des questions du type \og Quels articles utilisent SQuAD ? \fg{}, puis
rédige la réponse comme un RAG.
\begin{analogie}
Le RAG classique cherche des paragraphes qui ressemblent à la question. Le GraphRAG consulte
un arbre généalogique des connaissances : il suit les liens \og utilise \fg{},
\og évalué sur \fg{}, \og écrit par \fg{} pour trouver la réponse exacte.
\end{analogie}

\paragraph{Le topic modeling (BERTopic).}
Le \textbf{topic modeling} regroupe automatiquement des milliers de résumés en
\emph{thèmes}. BERTopic utilise les embeddings : UMAP réduit les vecteurs à quelques
dimensions, HDBSCAN forme des groupes de points denses, et chaque groupe reçoit un nom à
partir de ses mots-clés. On peut ensuite suivre l'évolution d'un thème mois par mois.
\begin{analogie}
C'est trier une pile de 20\,000 lettres sans les lire une par une, en formant des tas selon
leur sujet, puis en collant une étiquette sur chaque tas.
\end{analogie}

\paragraph{L'agent.}
Un \textbf{agent} est un LLM qui ne se contente pas de répondre : il \emph{choisit des
outils} (chercher des articles, interroger le graphe, lire un PDF\ldots), observe le
résultat, puis décide de l'étape suivante, jusqu'à pouvoir répondre. Cette boucle
\og réfléchir, agir, observer \fg{} s'appelle \textbf{ReAct}.
\begin{analogie}
Un agent est un stagiaire à qui l'on donne une boîte à outils et une mission : il décide
lui-même de chercher dans la bibliothèque, puis de consulter un annuaire, puis de rédiger.
\end{analogie}

\paragraph{LangGraph.}
\textbf{LangGraph} est une bibliothèque qui décrit un agent comme une \textbf{machine à
états} : des étapes (\og décider \fg{}, \og agir \fg{}, \og conclure \fg{}) reliées par des
flèches conditionnelles, avec un nombre maximal de tours pour éviter les boucles infinies.

% ======================================================================
\section{Le dataset arXiv}
% ======================================================================

\subsection{La source}

\textbf{arXiv} est le site où les chercheurs déposent leurs articles en accès libre.
L'université Cornell publie sur Kaggle un fichier qui contient les \textbf{métadonnées} de
tous les articles (titre, auteurs, résumé, catégories, dates\ldots) :
\url{https://www.kaggle.com/datasets/Cornell-University/arxiv}. C'est la \textbf{source de
données unique} du projet. Le fichier n'est jamais téléchargé automatiquement : il est
placé à la main dans \code{data/raw/arxiv-metadata-oai-snapshot.json}.

\subsection{Le format : JSON Lines}

Le fichier fait environ \textbf{5,6 Go}. Il est au format \textbf{JSON Lines} : chaque ligne
est un objet JSON complet qui décrit un article. Cela permet de le lire \textbf{ligne par
ligne} (en \emph{streaming}), sans jamais le charger entièrement en mémoire.

\begin{analogie}
Plutôt que de soulever un annuaire de 5 kilos d'un coup, on tourne les pages une à une, on
garde les fiches intéressantes et on oublie les autres.
\end{analogie}

\subsection{Un enregistrement expliqué}

\begin{lstlisting}
{"id": "0704.0001",
 "submitter": "Pavel Nadolsky",
 "authors": "C. Bal\'azs, E. L. Berger, P. M. Nadolsky, C.-P. Yuan",
 "title": "Calculation of prompt diphoton production cross sections ...",
 "comments": "37 pages, 15 figures; published version",
 "journal-ref": "Phys.Rev.D76:013009,2007",
 "doi": "10.1103/PhysRevD.76.013009",
 "categories": "hep-ph",
 "abstract": "  A fully differential calculation in perturbative ...",
 "versions": [{"version": "v1", "created": "Mon, 2 Apr 2007 19:18:42 GMT"}],
 "update_date": "2008-11-26",
 "authors_parsed": [["Balazs", "C.", ""], ["Berger", "E. L.", ""]]}
\end{lstlisting}

\begin{longtable}{>{\raggedright\arraybackslash}p{3.1cm}>{\raggedright\arraybackslash}p{11.8cm}}
\toprule
Champ & Signification et usage dans le projet \\
\midrule
\endhead
\code{id} & Identifiant arXiv (\code{0704.0001}, ou \code{2503.10677} pour les articles
récents). C'est lui qui apparaît dans les citations \code{[arXiv:ID]}. \\
\code{authors} & Auteurs en texte brut, avec parfois des accents écrits en LaTeX
(\code{Bal\textbackslash'azs}). \\
\code{authors\_parsed} & Les mêmes auteurs déjà découpés en [nom, prénom, suffixe] : le
projet les utilise de préférence. \\
\code{title}, \code{abstract} & Titre et résumé, avec des retours à la ligne et des espaces
multiples à nettoyer. Ce sont les textes que l'on indexe. \\
\code{categories} & Catégories séparées par des espaces, par exemple
\code{"cs.CL cs.LG"}. La première est la catégorie principale. \\
\code{comments} & Commentaire libre de l'auteur : \og 12 pages, 3 figures \fg{},
\og Accepted at ACL 2024 \fg{}, lien GitHub\ldots{} Source principale de l'enrichissement. \\
\code{journal-ref}, \code{doi} & Référence de publication et identifiant DOI, souvent vides
(\code{null}). \\
\code{versions} & Historique des versions. La date de la \textbf{première} version
(\code{versions[0].created}) est la vraie date de publication. \\
\code{update\_date} & Date de la dernière mise à jour : \textbf{pas} la date de
publication. \\
\bottomrule
\end{longtable}

\subsection{Le filtrage}

\begin{enumerate}
  \item \textbf{Pré-filtre rapide} : une ligne qui ne contient même pas les caractères
  \code{cs.CL} est ignorée \emph{sans être décodée}. Cela élimine 96\,\% des lignes
  (3\,092\,613 sur 3\,215\,335) et rend la lecture très rapide (environ 1 min 15 s).
  \item \textbf{Catégorie} : on garde un article si \emph{au moins une} de ses catégories
  est \code{cs.CL} (122\,720 articles).
  \item \textbf{Date} : publication à partir du 1\textsuperscript{er} janvier 2023
  (81\,503 articles).
  \item \textbf{Échantillon} : 20\,000 articles tirés au hasard de façon
  \textbf{reproductible} (graine 42) par \emph{échantillonnage par réservoir}, qui tire au
  sort en une seule lecture et avec une mémoire limitée.
\end{enumerate}

Répartition obtenue : 3\,304 articles en 2023, 5\,063 en 2024, 5\,916 en 2025 et 5\,717
en 2026. Tous ces paramètres se changent dans \code{config.yaml}.

\subsection{Le nettoyage}

\begin{itemize}
  \item suppression des retours à la ligne et des espaces multiples dans les titres et
  résumés ;
  \item conversion des accents LaTeX en vrais caractères (\code{Bal\textbackslash'azs}
  devient \og Balázs \fg{}) ;
  \item reconstruction des noms d'auteurs \og Prénom Nom \fg{} depuis
  \code{authors\_parsed} ;
  \item lecture de la date de publication au format RFC~2822
  (\og Mon, 2 Apr 2007 19:18:42 GMT \fg{}) ;
  \item écriture du résultat au format \textbf{Parquet} dans
  \code{data/processed/papers.parquet}.
\end{itemize}

\subsection{L'enrichissement}

Le fichier \code{src/rpe/data/enrich\_comments.py} lit les champs \code{comments},
\code{journal-ref} et \code{abstract} avec des \textbf{expressions régulières} (des motifs
de recherche dans le texte) et ajoute des colonnes utiles pour filtrer.

\begin{center}
\begin{tabular}{>{\raggedright\arraybackslash}p{2.6cm}>{\raggedright\arraybackslash}p{5.6cm}>{\raggedright\arraybackslash}p{6.2cm}}
\toprule
Colonne & Exemple de texte & Résultat \\
\midrule
\code{venue}, \code{venue\_year} & \og Accepted at ACL 2024 \fg{} & ACL, 2024 \\
\code{venue} & \og Findings of EMNLP 2023 \fg{} & EMNLP Findings \\
\code{is\_accepted} & \og Under review at ICLR 2025 \fg{} & venue ICLR, mais non accepté \\
\code{github\_url}, \code{has\_code} & \og Code: github.com/x/y \fg{} & lien GitHub, code
disponible \\
\code{n\_pages}, \code{n\_figures} & \og 8+2 pages, 3 figures \fg{} & 10 pages, 3 figures \\
\bottomrule
\end{tabular}
\end{center}

27 conférences et revues sont reconnues (ACL, EMNLP, NAACL, NeurIPS, ICML, ICLR\ldots). Sur
le corpus : 16,8\,\% des articles ont un lien vers du code, 36,3\,\% sont acceptés, 28,7\,\%
ont une venue identifiée.

\subsection{Les limites du dataset}

\begin{itemize}
  \item Ce sont des \textbf{métadonnées} : on n'a que le résumé, pas l'article entier (sauf
  pour les PDF téléchargés à part).
  \item Les résumés sont en \textbf{anglais} ; les questions marchent mieux en anglais.
  \item Le corpus est un \textbf{échantillon} de 20\,000 articles d'une seule catégorie :
  un article pertinent peut ne pas y être.
  \item Les commentaires sont libres : les règles ne couvrent pas toutes les formulations ni
  toutes les conférences (FAccT ou WMT, par exemple, ne sont pas reconnues).
  \item Le fichier est un \textbf{instantané} : il faut le retélécharger pour avoir les
  derniers articles.
\end{itemize}

% ======================================================================
\section{L'architecture du projet}
% ======================================================================

\subsection{Vue d'ensemble}

\begin{center}
\begin{tikzpicture}[node distance=0.55cm and 0.75cm]
  \node[boite,fill=donnees] (raw) {Dump arXiv\\JSONL 5,6 Go};
  \node[boite,fill=traitement,below=of raw] (prep) {Préparation\\filtre, nettoyage};
  \node[boite,fill=stockage,below=of prep] (pq) {papers.parquet\\20\,000 articles};
  \node[boite,fill=stockage,right=of raw] (chroma) {ChromaDB\\embeddings};
  \node[boite,fill=stockage,right=of prep] (bm) {BM25\\mots-clés};
  \node[boite,fill=stockage,right=of pq] (pdf) {PDF arXiv\\par section};
  \node[boite,fill=traitement,right=of chroma] (hyb) {Recherche\\hybride RRF};
  \node[boite,fill=llm,right=of bm] (rag) {RAG\\+ vérification};
  \node[boite,fill=traitement,right=of pdf] (cmp) {Comparaison\\+ rapport PDF};
  \node[boite,fill=interface,right=of rag,text width=2.3cm,minimum height=2.6cm] (ui)
       {Streamlit\\Chat\\Recherche\\Historique};
  \node[boite,fill=llm,above=0.5cm of ui,text width=2.3cm] (llm) {LLM\\Ollama, Groq\ldots};
  \draw[fleche] (raw) -- (prep);
  \draw[fleche] (prep) -- (pq);
  \draw[fleche] (pq.east) -- (chroma.west);
  \draw[fleche] (pq.east) -- (bm.west);
  \draw[fleche] (chroma) -- (hyb);
  \draw[fleche] (bm) -- (hyb);
  \draw[fleche] (hyb) -- (rag);
  \draw[pointille] (pdf) -- (rag);
  \draw[fleche] (rag) -- (cmp);
  \draw[fleche] (hyb.east) -- (ui);
  \draw[fleche] (rag) -- (ui);
  \draw[fleche] (cmp.east) -- (ui);
  \draw[pointille] (llm) -- (rag);
\end{tikzpicture}
\end{center}

\textbf{Légende} : en bleu les données brutes, en jaune les fichiers stockés sur disque, en
vert les traitements, en orange ce qui fait appel au LLM, en violet l'interface.

\begin{attention}
Chaque étape produit des \textbf{fichiers} que l'étape suivante relit. On peut donc
relancer une étape sans tout refaire, et reprendre un calcul interrompu. Un mode
\code{--dev} travaille sur 1\,000 articles, dans des fichiers séparés, pour tester vite.
\end{attention}

\subsection{Le trajet d'une question}

\begin{center}
\begin{tikzpicture}[node distance=0.5cm and 0.6cm]
  \node[boite,fill=interface] (q) {Question\\de l'utilisatrice};
  \node[boite,fill=stockage,right=of q,yshift=0.8cm] (b) {BM25\\100 candidats};
  \node[boite,fill=stockage,right=of q,yshift=-0.8cm] (d) {Embeddings\\100 candidats};
  \node[boite,fill=traitement,right=of b,yshift=-0.8cm] (f) {Fusion RRF\\top 5};
  \node[boite,fill=llm,below=1.3cm of f] (l) {LLM\\réponse citée};
  \node[boite,fill=traitement,left=of l] (v) {Vérification\\des citations};
  \node[boite,fill=interface,left=of v] (r) {Réponse +\\comparaison + PDF};
  \draw[fleche] (q) -- (b);
  \draw[fleche] (q) -- (d);
  \draw[fleche] (b) -- (f);
  \draw[fleche] (d) -- (f);
  \draw[fleche] (f) -- node[right,font=\footnotesize]{contexte numéroté} (l);
  \draw[fleche] (l) -- (v);
  \draw[fleche] (v) -- (r);
\end{tikzpicture}
\end{center}

\subsection{Les briques, une par une}

\subsubsection{La configuration}
\textbf{Rôle} : centraliser les réglages. \textbf{Fichiers} : \code{config.yaml} (catégories,
dates, tailles, chemins, seuils), \code{.env} (adresse et nom du LLM, clé d'API, jamais
publiée sur Git), \code{src/rpe/config.py} (lecture et vérification de ces réglages).
\textbf{Exemple} : passer de \code{cs.CL} à \code{cs.*} dans \code{config.yaml} suffit à
élargir le corpus à toute l'informatique.

\subsubsection{La préparation des données}
\textbf{Rôle} : transformer le fichier brut en un tableau propre. \textbf{Fichiers} :
\code{data/load.py} (lecture ligne par ligne, comptage des lignes invalides),
\code{data/clean.py} (filtres, nettoyage, échantillonnage), \code{data/enrich\_comments.py}
(conférence, code, pages), \code{data/prepare.py} (enchaînement et statistiques),
\code{data/store.py} (accès à un article par son identifiant).
\textbf{Commande} : \code{python -m rpe.cli prepare-data}.

\subsubsection{Les index et la recherche}
\textbf{Rôle} : retrouver vite les bons articles. \textbf{Fichiers} : \code{index/embed.py}
(calcul des embeddings par paquets de 2\,048 articles, sauvegardés au fur et à mesure),
\code{index/vector\_store.py} (ChromaDB), \code{index/bm25.py} (BM25 et découpage des mots),
\code{index/filters.py} (filtres année, catégorie, conférence, code), \code{index/hybrid.py}
(fusion RRF et extrait le plus pertinent du résumé), \code{index/build.py} (construction).
\textbf{Commandes} : \code{build-index}, puis
\code{search "retrieval-augmented generation" --year-min 2024 --has-code}.

\subsubsection{La couche LLM}
\textbf{Rôle} : un point d'entrée unique vers le modèle de langage. \textbf{Fichiers} :
\code{llm/client.py}, \code{llm/prompts.py}, \code{llm/schemas.py}.
\textbf{Fonctionnement} : avant chaque appel, le client regarde dans un \textbf{cache}
disque si la même question a déjà été posée (aucun appel en double). Sinon il appelle le
LLM avec un délai maximal ; en cas d'erreur réseau, il réessaie après 2, 4, puis 8
secondes. Pour une sortie structurée, il vérifie le JSON avec Pydantic et, s'il est faux,
redemande une correction. Il compte les appels, les réussites du cache et les erreurs.

\subsubsection{Le RAG et la vérification des citations}
\textbf{Rôle} : répondre en citant. \textbf{Fichiers} : \code{rag/qa.py} et
\code{rag/citations.py}. \textbf{Fonctionnement} : les 5 articles retenus sont numérotés
dans le prompt (\og [1] arXiv:2503.10677 | titre (année, conférence) \fg{} suivi du
résumé). Le LLM doit citer chaque affirmation ou répondre \code{INSUFFICIENT\_CONTEXT}.
Ensuite, \code{validate\_citations()} retire toute citation absente de ces 5 articles. Si
aucune citation valide ne reste, la réponse devient \og Je ne sais pas \fg{}.

\subsubsection{La comparaison des articles cités}
\textbf{Rôle} : décrire et comparer les articles cités. \textbf{Fichier} :
\code{rag/compare.py}. \textbf{Fonctionnement} : pour chaque article cité (5 au maximum),
on prend le texte du PDF s'il a été téléchargé, sinon le résumé. \textbf{Un seul appel}
groupé au LLM extrait problème, méthode, jeux de données, métriques, résultats et limites ;
le résultat est mis en cache par article. Une information absente devient
\og Non mentionné dans le résumé \fg{}. La pertinence (élevée, moyenne, faible) se calcule
sans LLM : c'est la part des mots de la question présents dans le titre et le résumé.
Un second appel rédige une synthèse comparative de 4 à 6 phrases, elle aussi vérifiée.

\subsubsection{Le texte intégral (PDF)}
\textbf{Rôle} : aller au-delà du résumé. \textbf{Fichiers} : \code{pdf/download.py}
(jusqu'à 300 PDF, 3 secondes entre deux requêtes pour respecter les serveurs d'arXiv),
\code{pdf/parse.py} (extraction du texte et découpage en sections : Introduction, Method,
Experiments, Results, Conclusion\ldots), \code{pdf/chunks.py} (passages de 1\,500 caractères
avec chevauchement, indexés dans ChromaDB), \code{pdf/reader.py}. \textbf{Exemple} : dans
le Chat, une question sur un seul article peut citer \code{[arXiv:ID \S Method]}.
\textbf{Commandes} : \code{fetch-pdfs} puis \code{index-pdfs}.

\subsubsection{Les rapports PDF}
\textbf{Rôle} : produire un document partageable. \textbf{Fichiers} :
\code{reports/pdf\_report.py} (mise en page avec ReportLab : page de titre, réponse avec
liens cliquables, fiche par article, tableau comparatif en paysage, synthèse, références,
pied de page) et \code{reports/chat\_report.py}. Le rapport réutilise la comparaison déjà
affichée : aucun nouvel appel au LLM.

\subsubsection{L'interface}
\textbf{Rôle} : rendre tout cela utilisable sans ligne de commande. \textbf{Fichiers} :
\code{app/streamlit\_app.py} (menu), \code{app/pages/chat.py}, \code{app/pages/search.py},
\code{ui/common.py} (cartes d'articles, liens de citation), \code{ui/conversations.py}
(historique : un fichier JSON par conversation dans \code{data/conversations/}),
\code{ui/compare\_view.py}, \code{ui/report\_view.py}, et \code{resources.py} qui charge une
seule fois les index et le LLM. \textbf{Commande} : \code{streamlit run app/streamlit\_app.py}.

\subsubsection{La ligne de commande et les tests}
\code{src/rpe/cli.py} propose six commandes : \code{prepare-data}, \code{build-index},
\code{search}, \code{ask} (avec \code{--pdf} pour le rapport), \code{fetch-pdfs},
\code{index-pdfs}. Le dossier \code{tests/} contient 114 tests automatisés qui utilisent un
\og faux LLM \fg{} aux réponses prévues : ils sont rapides, gratuits et reproductibles.

% ======================================================================
\section{Pourquoi ces outils ?}
% ======================================================================

Pour chaque choix : la raison, l'alternative principale et l'inconvénient accepté.

\begin{longtable}{>{\raggedright\arraybackslash}p{2.6cm}>{\raggedright\arraybackslash}p{4.6cm}>{\raggedright\arraybackslash}p{3.6cm}>{\raggedright\arraybackslash}p{3.8cm}}
\toprule
Outil & Pourquoi lui & Alternative & Inconvénient \\
\midrule
\endhead
Python 3.12 & Langage de référence de l'IA, toutes les bibliothèques existent &
Java, JavaScript & Plus lent que du C, sans importance ici \\
orjson & Décode le JSON plusieurs fois plus vite que le module standard, crucial pour
3,2 millions de lignes & \code{json} standard, ujson & Une dépendance de plus \\
pandas + Parquet & Tableaux faciles à manipuler ; Parquet est compressé, typé et rapide
à relire & CSV, polars & Parquet n'est pas lisible dans un éditeur de texte \\
bge-small-en-v1.5 (sentence-transformers) & Petit (384 dimensions), gratuit, local, très
bon en recherche & Embeddings payants d'OpenAI, bge-large & Anglais seulement ; moins précis
qu'un grand modèle \\
ChromaDB & Base vectorielle persistante, simple, avec filtres sur les métadonnées, sans
serveur & FAISS, Qdrant, pgvector & Moins performante à très grande échelle \\
rank-bm25 & Implémentation simple et sûre de BM25 & bm25s, Elasticsearch & Un peu lent sur
des millions de documents \\
Fusion RRF ($k=60$) & Combine des classements sans avoir à comparer des scores d'échelles
différentes ; standard et robuste & Moyenne pondérée des scores & Ignore l'écart de score
entre deux rangs \\
Client compatible OpenAI & Le même code parle à Ollama, Groq, Gemini ou OpenAI & SDK propre
à chaque fournisseur, LangChain & Les fonctions propres à un fournisseur ne sont pas
exploitées \\
Ollama + qwen2.5:3b & Local, gratuit, sans clé, données privées & Groq, Gemini (hébergés) &
Lent sur CPU (1 à 4 min), moins précis \\
Pydantic & Vérifie automatiquement les sorties du LLM & Vérifications écrites à la main &
Le LLM doit respecter un schéma \\
PyMuPDF & Extraction de texte rapide et fiable & pdfplumber, GROBID & Ne reconnaît pas la
structure logique : découpage par titres heuristique \\
httpx & Téléchargement HTTP simple, déjà installé avec le client OpenAI & requests & Aucun \\
ReportLab & PDF en Python, sans logiciel externe, avec accents (police DejaVuSans) &
WeasyPrint, LaTeX & Mise en page plus verbeuse à coder \\
Streamlit & Application web en Python pur, idéale pour un prototype & Gradio, React &
Moins personnalisable qu'une vraie application web \\
Typer & Ligne de commande lisible, générée à partir des fonctions & argparse, click &
Aucun notable \\
pytest + ruff & Tests et contrôle de qualité standards en Python & unittest, flake8 &
Aucun notable \\
\bottomrule
\end{longtable}

% ======================================================================
\section{Les garde-fous contre les hallucinations}
\label{sec:gardefous}
% ======================================================================

Le principe : \textbf{on ne croit jamais le LLM sur parole}, le programme vérifie.

\begin{longtable}{>{\raggedright\arraybackslash}p{5.2cm}>{\raggedright\arraybackslash}p{9.6cm}}
\toprule
Risque & Garde-fou \\
\midrule
\endhead
Citer un article inexistant ou absent du contexte & \code{validate\_citations()} retire
toute citation \code{[arXiv:ID]} qui ne fait pas partie des 5 articles récupérés. \\
Citer une section inventée (\code{\S 3.3}) & La section est retirée si elle n'existe pas
dans le contexte ; l'identifiant valide est conservé. \\
Répondre sans source & Si aucune citation valide ne reste, la réponse devient
\og Je ne sais pas \fg{}. Le LLM peut aussi répondre \code{INSUFFICIENT\_CONTEXT}. \\
Inventer un résultat chiffré dans la comparaison & Un résultat contenant un nombre absent du
texte source est écarté. \\
Inventer une information manquante & Les cases vides deviennent
\og Non mentionné dans le résumé \fg{}. \\
Références fausses dans le rapport & La liste des références est produite par le
programme à partir des seuls articles cités et vérifiés, jamais par le LLM. \\
JSON mal formé & Validation Pydantic, puis nouvelle demande au LLM avec le message
d'erreur. \\
Réponses instables & Température 0 (le LLM choisit toujours le mot le plus probable) et
cache : la même question donne la même réponse. \\
\bottomrule
\end{longtable}

\begin{attention}
Ces garde-fous garantissent que \textbf{les citations existent dans le contexte}. Ils ne
garantissent pas que chaque phrase résume parfaitement l'article cité : la lecture critique
reste nécessaire. C'est pourquoi le contexte récupéré est toujours affiché.
\end{attention}

% ======================================================================
\section{Limites et perspectives}
% ======================================================================

\subsection{Limites}
\begin{itemize}
  \item \textbf{Vitesse} : le LLM local de 3 milliards de paramètres répond en 1 à 4
  minutes sur un ordinateur sans carte graphique.
  \item \textbf{Qualité du LLM} : un petit modèle comprend moins bien les consignes ; il
  produit parfois des citations mal placées, que la vérification retire.
  \item \textbf{Couverture} : une seule catégorie, un échantillon de 20\,000 articles, des
  résumés en anglais.
  \item \textbf{Règles d'enrichissement} : certaines conférences et formulations ne sont
  pas reconnues.
  \item \textbf{PDF} : le découpage en sections repose sur les titres habituels et peut se
  tromper sur une mise en page atypique.
  \item \textbf{Évaluation} : la version livrée ne contient pas de mesure chiffrée de la
  qualité de recherche ou des réponses.
  \item \textbf{Périmètre} : les modules d'analyse du cahier des charges (thèmes, graphe,
  agent, évaluation) ont été retirés de la version livrée.
\end{itemize}

\subsection{Perspectives}
\begin{itemize}
  \item Utiliser un LLM hébergé (Groq, Gemini) : réponses en quelques secondes, sans
  changer le code.
  \item Ajouter un \textbf{re-classement} (cross-encoder) après la fusion RRF pour
  améliorer l'ordre des 5 articles retenus.
  \item \textbf{Mesurer} la qualité : Recall@k, MRR, nDCG pour la recherche ; taux de
  citations valides et fidélité pour les réponses, sur un jeu annoté à la main.
  \item Réintégrer l'analyse du domaine : thèmes et tendances, graphe de connaissances
  (GraphRAG), agent LangGraph.
  \item Élargir le corpus à d'autres catégories et le mettre à jour automatiquement avec
  l'API arXiv.
\end{itemize}

% ======================================================================
\section{Questions probables du jury}
% ======================================================================

\newcommand{\question}[1]{\paragraph{Q : #1}\mbox{}\\}

\question{Pourquoi ne pas simplement poser la question à ChatGPT ?}
Parce qu'un LLM seul répond de mémoire et peut inventer des références. Ici, il répond
uniquement à partir d'articles réellement retrouvés, et le programme vérifie chaque
citation. On sait toujours d'où vient une information.

\question{Qu'est-ce que le RAG, en une phrase ?}
Chercher d'abord les bons documents, puis demander au LLM de répondre uniquement avec eux,
en les citant.

\question{Pourquoi deux moteurs de recherche ?}
BM25 trouve les mots exacts (sigles, noms de jeux de données comme SQuAD) ; les embeddings
trouvent le sens (synonymes, reformulations). Leur combinaison est meilleure que chacun
seul, surtout sur des requêtes courtes et techniques.

\question{Comment fonctionne la fusion RRF ?}
Chaque article reçoit $\frac{1}{60+\text{rang}}$ pour chaque moteur, et on additionne. Un
article bien classé par les deux moteurs passe devant un article très bien classé par un
seul. On garde ensuite les 5 premiers.

\question{Pourquoi garder 5 articles et pas plus ?}
Plus d'articles allongent le prompt, ralentissent un petit LLM local et diluent
l'information. Cinq articles suffisent pour une réponse ciblée. La valeur se règle dans
\code{config.yaml}.

\question{Comment êtes-vous sûres que les citations ne sont pas inventées ?}
La fonction \code{validate\_citations()} compare chaque identifiant cité à la liste des
articles récupérés et retire ceux qui n'y sont pas. Sans citation valide, la réponse devient
\og Je ne sais pas \fg{}. Cela est vérifié par des tests automatisés.

\question{Que se passe-t-il si la question sort du sujet du corpus ?}
Les articles récupérés ne contiennent pas la réponse : le LLM doit répondre
\code{INSUFFICIENT\_CONTEXT}, ou ses citations sont retirées, et l'application affiche
\og Je ne sais pas \fg{}.

\question{Pourquoi un modèle local plutôt qu'un modèle en ligne ?}
Gratuit, sans clé d'API, et les données restent sur la machine. L'inconvénient est la
lenteur. Le code est conçu pour changer de fournisseur en modifiant seulement le fichier
\code{.env}.

\question{Pourquoi seulement 20\,000 articles ?}
Pour que les embeddings et les index se calculent en un temps raisonnable sur un ordinateur
portable. Le tirage est reproductible (graine 42) et la taille se règle dans
\code{config.yaml}.

\question{Comment lisez-vous un fichier de 5,6 Go ?}
Ligne par ligne, sans jamais le charger en entier. Un pré-filtre ignore 96\,\% des lignes
sans même les décoder. La lecture complète prend environ une minute et quart.

\question{Pourquoi la date de publication vient-elle de \code{versions} ?}
Parce que \code{update\_date} est la date de la dernière modification, qui peut être des
années après la publication. La première version donne la vraie date.

\question{Comment détectez-vous la conférence d'un article ?}
Avec des expressions régulières sur le commentaire de l'auteur et la référence de
publication (\og Accepted at ACL 2024 \fg{}, \og To appear in NeurIPS'23 \fg{}), en
repérant aussi les formulations négatives comme \og under review \fg{}.

\question{Que fait la comparaison des articles ?}
Pour les articles cités, elle extrait en un seul appel au LLM la méthode, les jeux de
données, les résultats et les limites, puis rédige une synthèse. Ce qui n'est pas écrit
dans le texte est marqué \og Non mentionné \fg{} et les chiffres absents du texte sont
écartés.

\question{Qu'est-ce qu'un embedding ?}
Une liste de nombres qui représente le sens d'un texte. Deux textes de sens proche ont des
listes proches, ce qui permet de chercher par le sens et non par les mots exacts.

\question{Qu'est-ce que le GraphRAG et pourquoi n'est-il pas dans la version livrée ?}
C'est un RAG qui interroge un graphe de connaissances (articles, tâches, méthodes, jeux de
données reliés entre eux). Il a été développé et testé, puis retiré de la version livrée
pour la concentrer sur la recherche et le RAG ; il reste dans l'historique Git.

\question{Qu'est-ce qu'un agent, et LangGraph ?}
Un agent est un LLM qui choisit lui-même des outils et des étapes pour accomplir une tâche.
LangGraph décrit cet agent comme une machine à états avec un nombre maximal de tours. Comme
le GraphRAG, l'agent a été retiré de la version livrée.

\question{Comment avez-vous testé le projet ?}
114 tests automatisés vérifient la lecture des données, l'enrichissement, la recherche, la
fusion, les citations, la comparaison, le rapport PDF et l'interface. Ils utilisent un faux
LLM aux réponses prévues, donc ils sont rapides et reproductibles. La qualité du code est
contrôlée par ruff.

\question{Quelle est la principale limite ?}
La vitesse et la précision du petit LLM local. La solution la plus simple est d'utiliser un
LLM hébergé, ce que le code permet sans modification.

% ======================================================================
\appendix
\section{Commandes pour lancer le projet}
% ======================================================================

\begin{lstlisting}
py -3.12 -m venv .venv
.venv\Scripts\activate
pip install -r requirements.txt
pip install -e . --no-deps
copy .env.example .env               (puis choisir le LLM)

python -m rpe.cli prepare-data       (corpus ; --dev pour 1 000 articles)
python -m rpe.cli build-index        (embeddings, ChromaDB, BM25)
python -m rpe.cli fetch-pdfs         (optionnel : PDF arXiv)
python -m rpe.cli index-pdfs
streamlit run app/streamlit_app.py   (http://localhost:8501)

python -m rpe.cli ask "How can retrieval reduce hallucinations in LLMs?" --lang fr --pdf rapport.pdf
pytest
ruff check .
\end{lstlisting}

\end{document}
